\documentclass[10pt,a4paper]{amsart}
\usepackage{amsmath,amssymb,amsthm,graphicx}
% Retained preamble for the Stacks Project chapter {CH} (shared project preamble).
% The chapter file pinned-original-source.tex only carries \input{{preamble}} on its first line,
% so no preamble text lives inside the chapter itself.  The Stacks Project uses one single shared
% preamble for every chapter: the snapshot stages it three times (sources/stacks-chow/preamble.tex,
% sources/stacks-etale-cohomology/preamble.tex and focused-seeds/AGP/sources/stacks-snapshot-20261011/
% preamble.tex) and those copies are identical in size.  The block below is copied mechanically,
% line for line, from sources/stacks-chow/preamble.tex lines 14-18 and 154-217: the Xy-pic setup and
% the native theorem declarations together with the mathematical macros.  Nothing was added, removed
% or redefined; the wrapper class and the amsmath/amssymb/amsthm/graphicx base packages are supplied
% by the assembler.
% For commutative diagrams we use Xy-pic
% Chapter: stacks-spaces-more-morphisms.tex -- More on Morphisms of Spaces.
\usepackage[all]{xy}

% We use 2cell for 2-commutative diagrams.
\xyoption{2cell}

% Theorem environments.
%
\theoremstyle{plain}
\newtheorem{theorem}[subsection]{Theorem}
\newtheorem{proposition}[subsection]{Proposition}
\newtheorem{lemma}[subsection]{Lemma}

\theoremstyle{definition}
\newtheorem{definition}[subsection]{Definition}
\newtheorem{example}[subsection]{Example}
\newtheorem{exercise}[subsection]{Exercise}
\newtheorem{situation}[subsection]{Situation}

\theoremstyle{remark}
\newtheorem{remark}[subsection]{Remark}
\newtheorem{remarks}[subsection]{Remarks}

\numberwithin{equation}{subsection}

% Macros
%
\def\lim{\mathop{\mathrm{lim}}\nolimits}
\def\colim{\mathop{\mathrm{colim}}\nolimits}
\def\Spec{\mathop{\mathrm{Spec}}}
\def\Hom{\mathop{\mathrm{Hom}}\nolimits}
\def\Ext{\mathop{\mathrm{Ext}}\nolimits}
\def\SheafHom{\mathop{\mathcal{H}\!\mathit{om}}\nolimits}
\def\SheafExt{\mathop{\mathcal{E}\!\mathit{xt}}\nolimits}
\def\Sch{\mathit{Sch}}
\def\Mor{\mathop{\mathrm{Mor}}\nolimits}
\def\Ob{\mathop{\mathrm{Ob}}\nolimits}
\def\Sh{\mathop{\mathit{Sh}}\nolimits}
\def\NL{\mathop{N\!L}\nolimits}
\def\CH{\mathop{\mathrm{CH}}\nolimits}
\def\proetale{{pro\text{-}\acute{e}tale}}
\def\etale{{\acute{e}tale}}
\def\QCoh{\mathit{QCoh}}
\def\Ker{\mathop{\mathrm{Ker}}}
\def\Im{\mathop{\mathrm{Im}}}
\def\Coker{\mathop{\mathrm{Coker}}}
\def\Coim{\mathop{\mathrm{Coim}}}

% Boxtimes
%
\DeclareMathSymbol{\boxtimes}{\mathbin}{AMSa}{"02}

%
% Macros for moduli stacks/spaces
%
\def\QCohstack{\mathcal{QC}\!\mathit{oh}}
\def\Cohstack{\mathcal{C}\!\mathit{oh}}
\def\Spacesstack{\mathcal{S}\!\mathit{paces}}
\def\Quotfunctor{\mathrm{Quot}}
\def\Hilbfunctor{\mathrm{Hilb}}
\def\Curvesstack{\mathcal{C}\!\mathit{urves}}
\def\Polarizedstack{\mathcal{P}\!\mathit{olarized}}
\def\Complexesstack{\mathcal{C}\!\mathit{omplexes}}
% \Pic is the operator that assigns to X its picard group, usage \Pic(X)
% \Picardstack_{X/B} denotes the Picard stack of X over B
% \Picardfunctor_{X/B} denotes the Picard functor of X over B
\def\Pic{\mathop{\mathrm{Pic}}\nolimits}
\def\Picardstack{\mathcal{P}\!\mathit{ic}}
\def\Picardfunctor{\mathrm{Pic}}
\def\Deformationcategory{\mathcal{D}\!\mathit{ef}}
\newenvironment{reference}{}{}
\newenvironment{slogan}{}{}
\newenvironment{history}{}{}
\usepackage{cleveref}
\crefname{originallemma}{original lemma}{original lemmas}
\Crefname{originallemma}{Original lemma}{Original lemmas}
\makeatletter
\newlabel{definition-formally-smooth}{{definition-formally-smooth}{source}{}{sourcepointer-1}{}}
\newlabel{definition-formally-smooth@cref}{{[originallemma][0][]definition-formally-smooth}{[1][0][]source}{}{}{}}
\makeatother
\crefname{originallemma}{original lemma}{original lemmas}
\Crefname{originallemma}{Original lemma}{Original lemmas}
\makeatletter
\newlabel{lemma-action-sheaf}{{lemma-action-sheaf}{source}{}{sourcepointer-2}{}}
\newlabel{lemma-action-sheaf@cref}{{[originallemma][0][]lemma-action-sheaf}{[1][0][]source}{}{}{}}
\makeatother
\crefname{originallemma}{original lemma}{original lemmas}
\Crefname{originallemma}{Original lemma}{Original lemmas}
\makeatletter
\newlabel{lemma-composition-formally-smooth-etale-unramified}{{lemma-composition-formally-smooth-etale-unramified}{source}{}{sourcepointer-3}{}}
\newlabel{lemma-composition-formally-smooth-etale-unramified@cref}{{[originallemma][0][]lemma-composition-formally-smooth-etale-unramified}{[1][0][]source}{}{}{}}
\makeatother
\crefname{originallemma}{original lemma}{original lemmas}
\Crefname{originallemma}{Original lemma}{Original lemmas}
\makeatletter
\newlabel{lemma-etale-on-top}{{lemma-etale-on-top}{source}{}{sourcepointer-4}{}}
\newlabel{lemma-etale-on-top@cref}{{[originallemma][0][]lemma-etale-on-top}{[1][0][]source}{}{}{}}
\makeatother
\crefname{originallemma}{original lemma}{original lemmas}
\Crefname{originallemma}{Original lemma}{Original lemmas}
\makeatletter
\newlabel{lemma-formally-smooth-sheaf-differentials}{{lemma-formally-smooth-sheaf-differentials}{source}{}{sourcepointer-5}{}}
\newlabel{lemma-formally-smooth-sheaf-differentials@cref}{{[originallemma][0][]lemma-formally-smooth-sheaf-differentials}{[1][0][]source}{}{}{}}
\makeatother
\crefname{originallemma}{original lemma}{original lemmas}
\Crefname{originallemma}{Original lemma}{Original lemmas}
\makeatletter
\newlabel{lemma-h1-is-zero}{{lemma-h1-is-zero}{source}{}{sourcepointer-6}{}}
\newlabel{lemma-h1-is-zero@cref}{{[originallemma][0][]lemma-h1-is-zero}{[1][0][]source}{}{}{}}
\makeatother
\crefname{originallemma}{original lemma}{original lemmas}
\Crefname{originallemma}{Original lemma}{Original lemmas}
\makeatletter
\newlabel{lemma-helper-formally-smooth}{{lemma-helper-formally-smooth}{source}{}{sourcepointer-7}{}}
\newlabel{lemma-helper-formally-smooth@cref}{{[originallemma][0][]lemma-helper-formally-smooth}{[1][0][]source}{}{}{}}
\makeatother
\crefname{originallemma}{original lemma}{original lemmas}
\Crefname{originallemma}{Original lemma}{Original lemmas}
\makeatletter
\newlabel{lemma-representable-property-formally-property}{{lemma-representable-property-formally-property}{source}{}{sourcepointer-8}{}}
\newlabel{lemma-representable-property-formally-property@cref}{{[originallemma][0][]lemma-representable-property-formally-property}{[1][0][]source}{}{}{}}
\makeatother
\crefname{originallemma}{original lemma}{original lemmas}
\Crefname{originallemma}{Original lemma}{Original lemmas}
\makeatletter
\newlabel{lemma-sheaf}{{lemma-sheaf}{source}{}{sourcepointer-9}{}}
\newlabel{lemma-sheaf@cref}{{[originallemma][0][]lemma-sheaf}{[1][0][]source}{}{}{}}
\makeatother
\crefname{originallemma}{original lemma}{original lemmas}
\Crefname{originallemma}{Original lemma}{Original lemmas}
\makeatletter
\newlabel{lemma-smooth-formally-smooth}{{lemma-smooth-formally-smooth}{source}{}{sourcepointer-10}{}}
\newlabel{lemma-smooth-formally-smooth@cref}{{[originallemma][0][]lemma-smooth-formally-smooth}{[1][0][]source}{}{}{}}
\makeatother
\crefname{originallemma}{original lemma}{original lemmas}
\Crefname{originallemma}{Original lemma}{Original lemmas}
\makeatletter
\newlabel{more-morphisms-definition-formally-unramified}{{more-morphisms-definition-formally-unramified}{source}{}{sourcepointer-11}{}}
\newlabel{more-morphisms-definition-formally-unramified@cref}{{[originallemma][0][]more-morphisms-definition-formally-unramified}{[1][0][]source}{}{}{}}
\makeatother
\crefname{originallemma}{original lemma}{original lemmas}
\Crefname{originallemma}{Original lemma}{Original lemmas}
\makeatletter
\newlabel{remark-special-case}{{remark-special-case}{source}{}{sourcepointer-12}{}}
\newlabel{remark-special-case@cref}{{[originallemma][0][]remark-special-case}{[1][0][]source}{}{}{}}
\makeatother
\crefname{originallemma}{original lemma}{original lemmas}
\Crefname{originallemma}{Original lemma}{Original lemmas}
\makeatletter
\newlabel{sites-cohomology-definition-torsor}{{sites-cohomology-definition-torsor}{source}{}{sourcepointer-13}{}}
\newlabel{sites-cohomology-definition-torsor@cref}{{[originallemma][0][]sites-cohomology-definition-torsor}{[1][0][]source}{}{}{}}
\makeatother
\crefname{originallemma}{original lemma}{original lemmas}
\Crefname{originallemma}{Original lemma}{Original lemmas}
\makeatletter
\newlabel{sites-cohomology-lemma-torsors-h1}{{sites-cohomology-lemma-torsors-h1}{source}{}{sourcepointer-14}{}}
\newlabel{sites-cohomology-lemma-torsors-h1@cref}{{[originallemma][0][]sites-cohomology-lemma-torsors-h1}{[1][0][]source}{}{}{}}
\makeatother
\crefname{originallemma}{original lemma}{original lemmas}
\Crefname{originallemma}{Original lemma}{Original lemmas}
\makeatletter
\newlabel{spaces-descent-lemma-locally-projective-descends}{{spaces-descent-lemma-locally-projective-descends}{source}{}{sourcepointer-15}{}}
\newlabel{spaces-descent-lemma-locally-projective-descends@cref}{{[originallemma][0][]spaces-descent-lemma-locally-projective-descends}{[1][0][]source}{}{}{}}
\makeatother
\crefname{originallemma}{original lemma}{original lemmas}
\Crefname{originallemma}{Original lemma}{Original lemmas}
\makeatletter
\newlabel{spaces-properties-lemma-locally-projective-pullback}{{spaces-properties-lemma-locally-projective-pullback}{source}{}{sourcepointer-16}{}}
\newlabel{spaces-properties-lemma-locally-projective-pullback@cref}{{[originallemma][0][]spaces-properties-lemma-locally-projective-pullback}{[1][0][]source}{}{}{}}
\makeatother
\title{The infinitesimal lifting criterion for formally smooth morphisms of algebraic spaces}
\author{The Stacks Project Authors}
\date{stacks-project:spaces-more-morphisms.tex:local-snapshot-20261011; GNU FDL 1.2 or later}
\begin{document}
\maketitle
\noindent\textit{Editorial note.} Statement span is lines 4836-4857 of sources/stacks-spaces-more-morphisms/pinned-original-source.tex; proof span is lines 4859-4935, the complete author proof (75 body lines).\par\medskip
\noindent\textit{Editorial note.} Navigation stubs: 16 cited-but-not-included labels are supplied as non-load-bearing pointers with their pinned tags where the snapshot tag index has them.\par\medskip
\noindent\textit{Editorial note.} extra\_packages is empty on purpose: the retained preamble already loads xy with options and a second declaration clashes in the assembler wrapper.\par\medskip
\noindent\textit{Editorial note.} Text is retained verbatim, including the author's own typos and placeholders.\par\medskip
\noindent\textit{Allowed standard prerequisites (editorial).}\begin{itemize}
\item Explicit prerequisite boundary. Formally smooth morphisms, the small etale site, and the cohomological obstruction to smoothness; the implication is retained rather than cited.
\end{itemize}\medskip
\ifdefined\hypertarget\hypertarget{sourcepointer-1}{}\fi
\noindent\textit{Source reference.} Original lemma definition-formally-smooth: sources/stacks-spaces-more-morphisms/pinned-original-source.tex, source lines 4491--4496. Cited by the retained statement or proof but not part of the collected spans; supplied only as a navigation pointer. The argument does not depend on its content.\par\medskip
\ifdefined\hypertarget\hypertarget{sourcepointer-2}{}\fi
\noindent\textit{Source reference.} Original lemma lemma-action-sheaf: sources/stacks-spaces-more-morphisms/pinned-original-source.tex, source lines 3827--3838. Cited by the retained statement or proof but not part of the collected spans; supplied only as a navigation pointer. The argument does not depend on its content.\par\medskip
\ifdefined\hypertarget\hypertarget{sourcepointer-3}{}\fi
\noindent\textit{Source reference.} Original lemma lemma-composition-formally-smooth-etale-unramified: sources/stacks-spaces-more-morphisms/pinned-original-source.tex, source lines 2486--2495. Cited by the retained statement or proof but not part of the collected spans; supplied only as a navigation pointer. The argument does not depend on its content.\par\medskip
\ifdefined\hypertarget\hypertarget{sourcepointer-4}{}\fi
\noindent\textit{Source reference.} Original lemma lemma-etale-on-top: sources/stacks-spaces-more-morphisms/pinned-original-source.tex, source lines 2567--2576. Cited by the retained statement or proof but not part of the collected spans; supplied only as a navigation pointer. The argument does not depend on its content.\par\medskip
\ifdefined\hypertarget\hypertarget{sourcepointer-5}{}\fi
\noindent\textit{Source reference.} Original lemma lemma-formally-smooth-sheaf-differentials: sources/stacks-spaces-more-morphisms/pinned-original-source.tex, source lines 4769--4773. Cited by the retained statement or proof but not part of the collected spans; supplied only as a navigation pointer. The argument does not depend on its content.\par\medskip
\ifdefined\hypertarget\hypertarget{sourcepointer-6}{}\fi
\noindent\textit{Source reference.} Original lemma lemma-h1-is-zero: sources/stacks-spaces-more-morphisms/pinned-original-source.tex, source lines 4803--4812. Cited by the retained statement or proof but not part of the collected spans; supplied only as a navigation pointer. The argument does not depend on its content.\par\medskip
\ifdefined\hypertarget\hypertarget{sourcepointer-7}{}\fi
\noindent\textit{Source reference.} Original lemma lemma-helper-formally-smooth: sources/stacks-spaces-more-morphisms/pinned-original-source.tex, source lines 4554--4565. Cited by the retained statement or proof but not part of the collected spans; supplied only as a navigation pointer. The argument does not depend on its content.\par\medskip
\ifdefined\hypertarget\hypertarget{sourcepointer-8}{}\fi
\noindent\textit{Source reference.} Original lemma lemma-representable-property-formally-property: sources/stacks-spaces-more-morphisms/pinned-original-source.tex, source lines 2529--2537. Cited by the retained statement or proof but not part of the collected spans; supplied only as a navigation pointer. The argument does not depend on its content.\par\medskip
\ifdefined\hypertarget\hypertarget{sourcepointer-9}{}\fi
\noindent\textit{Source reference.} Original lemma lemma-sheaf: sources/stacks-spaces-more-morphisms/pinned-original-source.tex, source lines 3791--3805. Cited by the retained statement or proof but not part of the collected spans; supplied only as a navigation pointer. The argument does not depend on its content.\par\medskip
\ifdefined\hypertarget\hypertarget{sourcepointer-10}{}\fi
\noindent\textit{Source reference.} Original lemma lemma-smooth-formally-smooth: sources/stacks-spaces-more-morphisms/pinned-original-source.tex, source lines 4587--4595. Cited by the retained statement or proof but not part of the collected spans; supplied only as a navigation pointer. The argument does not depend on its content.\par\medskip
\ifdefined\hypertarget\hypertarget{sourcepointer-11}{}\fi
\noindent\textit{Source reference.} Original lemma more-morphisms-definition-formally-unramified: sources/stacks-spaces-more-morphisms/pinned-original-source.tex, source lines 4935--4935. Cited by the retained statement or proof but not part of the collected spans; supplied only as a navigation pointer. The argument does not depend on its content.\par\medskip
\ifdefined\hypertarget\hypertarget{sourcepointer-12}{}\fi
\noindent\textit{Source reference.} Original lemma remark-special-case: sources/stacks-spaces-more-morphisms/pinned-original-source.tex, source lines 3848--3864. Cited by the retained statement or proof but not part of the collected spans; supplied only as a navigation pointer. The argument does not depend on its content.\par\medskip
\ifdefined\hypertarget\hypertarget{sourcepointer-13}{}\fi
\noindent\textit{Source reference.} Original lemma sites-cohomology-definition-torsor: sources/stacks-spaces-more-morphisms/pinned-original-source.tex, source lines 4935--4935. Cited by the retained statement or proof but not part of the collected spans; supplied only as a navigation pointer. The argument does not depend on its content.\par\medskip
\ifdefined\hypertarget\hypertarget{sourcepointer-14}{}\fi
\noindent\textit{Source reference.} Original lemma sites-cohomology-lemma-torsors-h1: sources/stacks-spaces-more-morphisms/pinned-original-source.tex, source lines 4935--4935. Cited by the retained statement or proof but not part of the collected spans; supplied only as a navigation pointer. The argument does not depend on its content.\par\medskip
\ifdefined\hypertarget\hypertarget{sourcepointer-15}{}\fi
\noindent\textit{Source reference.} Original lemma spaces-descent-lemma-locally-projective-descends: sources/stacks-spaces-more-morphisms/pinned-original-source.tex, source lines 4935--4935. Cited by the retained statement or proof but not part of the collected spans; supplied only as a navigation pointer. The argument does not depend on its content.\par\medskip
\ifdefined\hypertarget\hypertarget{sourcepointer-16}{}\fi
\noindent\textit{Source reference.} Original lemma spaces-properties-lemma-locally-projective-pullback: sources/stacks-spaces-more-morphisms/pinned-original-source.tex, source lines 4935--4935. Cited by the retained statement or proof but not part of the collected spans; supplied only as a navigation pointer. The argument does not depend on its content.\par\medskip
\begin{lemma}
\label{lemma-formally-smooth}
Let $S$ be a scheme.
Let $f : X \to Y$ be a morphism of algebraic spaces over
$S$. The following are equivalent:
\begin{enumerate}
\item $f$ is formally smooth,
\item for every diagram
$$
\xymatrix{
U \ar[d] \ar[r]_\psi & V \ar[d] \\
X \ar[r]^f & Y
}
$$
where $U$ and $V$ are schemes and the vertical arrows are \'etale
the morphism of schemes $\psi$ is formally smooth (as in
More on Morphisms,
Definition \ref{more-morphisms-definition-formally-unramified}), and
\item for one such diagram with surjective vertical arrows the morphism
$\psi$ is formally smooth.
\end{enumerate}
\end{lemma}

\begin{proof}
We have seen that (1) implies (2) and (3) in
Lemma \ref{lemma-helper-formally-smooth}.
Assume (3).
The proof that $f$ is formally smooth is entirely similar to
the proof of (1) $\Rightarrow$ (2) of
Lemma \ref{lemma-smooth-formally-smooth}.

\medskip\noindent
Consider a solid commutative diagram
$$
\xymatrix{
X \ar[d]_f & T \ar[d]^i \ar[l]^a \\
Y & T' \ar[l] \ar@{-->}[lu]
}
$$
as in Definition \ref{definition-formally-smooth}.
We will show the dotted arrow exists thereby
proving that $f$ is formally smooth.
Let $\mathcal{F}$ be the sheaf of sets on $(T')_{spaces, \etale}$ of
Lemma \ref{lemma-sheaf} as in the special case discussed in
Remark \ref{remark-special-case}.
Let
$$
\mathcal{H} =
\SheafHom_{\mathcal{O}_T}(a^*\Omega_{X/Y}, \mathcal{C}_{T/T'})
$$
be the sheaf of $\mathcal{O}_T$-modules on $T_{spaces, \etale}$
with action $\mathcal{H} \times \mathcal{F} \to \mathcal{F}$ as in
Lemma \ref{lemma-action-sheaf}.
The action $\mathcal{H} \times \mathcal{F} \to \mathcal{F}$
turns $\mathcal{F}$ into a pseudo $\mathcal{H}$-torsor, see
Cohomology on Sites, Definition \ref{sites-cohomology-definition-torsor}.
Our goal is to show that $\mathcal{F}$ is a trivial $\mathcal{H}$-torsor.
There are two steps: (I) To show that $\mathcal{F}$ is a torsor
we have to show that $\mathcal{F}$ has \'etale locally a
section. (II) To show that $\mathcal{F}$ is the trivial torsor
it suffices to show that $H^1(T_\etale, \mathcal{H}) = 0$, see
Cohomology on Sites, Lemma \ref{sites-cohomology-lemma-torsors-h1}.

\medskip\noindent
First we prove (I). To see this consider a diagram
(which exists because we are assuming (3))
$$
\xymatrix{
U \ar[d] \ar[r]_\psi & V \ar[d] \\
X \ar[r]^f & Y
}
$$
where $U$ and $V$ are schemes, the vertical arrows are \'etale and
surjective, and $\psi$ is formally smooth. By
Lemma \ref{lemma-representable-property-formally-property}
the morphism $V \to Y$ is formally \'etale. Thus by
Lemma \ref{lemma-composition-formally-smooth-etale-unramified}
the composition $U \to Y$ is formally smooth. Then (I) follows from
Lemma \ref{lemma-etale-on-top} part (4).

\medskip\noindent
Finally we prove (II). By
Lemma \ref{lemma-formally-smooth-sheaf-differentials}
we see that $\Omega_{U/V}$ locally projective.
Hence $\Omega_{X/Y}$ is locally projective, see
Descent on Spaces,
Lemma \ref{spaces-descent-lemma-locally-projective-descends}.
Hence $a^*\Omega_{X/Y}$ is locally projective, see
Properties of Spaces, Lemma
\ref{spaces-properties-lemma-locally-projective-pullback}.
Hence
$$
H^1(T_\etale, \mathcal{H}) =
H^1(T_\etale,
\SheafHom_{\mathcal{O}_T}(a^*\Omega_{X/Y}, \mathcal{C}_{T/T'}) = 0
$$
by
Lemma \ref{lemma-h1-is-zero}
as desired.
\end{proof}

\end{document}
